\section{Genealogía completa de la proyección excepcional: de Paley y
Hadamard a los sistemas de Witt, Leech y Moonshine}
\label{sec:genealogia-excepcional-completa-autonoma}
\index{doble proyección!genealogía excepcional}
\index{Golay!ternario y binario}
\index{Steiner!sistemas de Witt}
\index{Mathieu!grupos M12 y M24}
\index{Moonshine!genealogía}

La evaluación excepcional parte del mismo estado enriquecido que la
evaluación arquimediana. La primera conserva incidencia, orientación,
acarreo y frontera; la segunda publica el valor real de la sección. Esta
bifurcación hace coexistir, desde un antecedente común, una genealogía de
valores y una genealogía de simetrías.

Dos transformaciones lineales cumplen funciones sucesivas y tipadas. La
transformada de Hadamard organiza el estado dodecafásico firmado y sus
sectores ciclotómicos. La matriz de conferencia de Paley fija después el
calibre proyectivo de la rama excepcional. Su reducción ternaria
\[
 C_P C_P^{\mathsf T}=5I_6,
 \qquad
 A_W=C_P\bmod3,
 \qquad
 A_W^2=-I_6
\]
produce el endomorfismo cuadrático interno y el código gráfico
\[
 C_W=\{(q,qA_W):q\in\mathbb F_3^6\}
      =[12,6,6]_3.
\]
El enumerador
\[
 W_{C_W}(y)=1+264y^6+440y^9+24y^{12}
\]
expresa la métrica de Hamming de esta realización. Los 264 vectores de peso
seis, identificados por pares opuestos, determinan 132 hexadas y construyen
el primer sistema de Witt,
\[
 W_{12}=S(5,6,12).
\]
Las 495 estrellas de Witt proporcionan las involuciones de incidencia cuya
acción generada tiene imagen \(M_{12}\).

La segunda estructura de Steiner se obtiene en el código binario extendido
de Golay \(\mathcal G_{24}\). Los soportes de peso ocho forman
\[
 W_{24}=S(5,8,24),
 \qquad \operatorname{Aut}(W_{24})\cong M_{24}.
\]
Fijada una dodecada \(D\), las intersecciones de las octadas con \(D\) en
seis puntos reconstruyen \(W_{12}\); cada hexada admite una elevación única
a una octada. La pareja de sistemas
\[
 S(5,6,12)\longrightarrow S(5,8,24)
\]
es, por tanto, una prolongación de incidencia demostrada dentro del Golay
binario y conserva tanto la dodecada como su alineamiento.

Desde el código ternario común nace simultáneamente la rama reticular. El
pegado coordenado por coordenado sobre \(A_2^{12}\) construye el retículo de
Niemeier
\[
 C_W\longrightarrow N(A_2^{12}).
\]
El vector radial seleccionado define un vecino de índice tres par,
unimodular, de rango veinticuatro y sin raíces; la clasificación lo identifica
con el retículo de Leech:
\[
 N(A_2^{12})\xrightarrow{\;3\text{-vecino}\;}\Lambda_{24}.
\]
Sobre esta salida se construye el álgebra de operadores de vértice reticular
\(V_{\Lambda_{24}}\). La involución canónica de signo y su módulo torcido
producen el orbifold
\[
 V^\natural=V_{\Lambda_{24}}^+
 \oplus V_{\Lambda_{24}}^{T,+},
 \qquad
 \operatorname{Aut}(V^\natural)\cong\mathbb M.
\]
Su carácter graduado es el Hauptmodul normalizado
\[
 \operatorname{Tr}_{V^\natural}q^{L_0-1}
 =J(\tau)=j(\tau)-744
 =q^{-1}+196884q+21493760q^2+\cdots,
\]
y la teoría de Moonshine asigna a cada clase de conjugación
\([g]\subset\mathbb M\) la serie de McKay--Thompson
\(T_g(\tau)=\operatorname{Tr}_{V^\natural}
(g\,q^{L_0-1})\), dotada de las propiedades modulares y de replicabilidad
correspondientes.

El antecedente común de ambas prolongaciones conserva la cadena completa
\[
 w_6\longrightarrow w_{30}\longrightarrow R_{36}\longrightarrow G_9
 \longrightarrow x_{\rm term}\longrightarrow U_{12}^{\rm sgn}
 \longrightarrow K\longrightarrow B_K.
\]
Desde este estado terminal deben distinguirse la ruta de evaluación y la ruta
del calibre. El estado produce las palabras de incidencia; la matriz de Paley
define, por reducción ternaria, el operador que las evalúa. La genealogía
excepcional completa queda por tanto organizada como sigue:
\[
\begin{aligned}
C_P&\longmapsto A_W=C_P\bmod3,\\
x\in X_\infty^{\rm cal}
&\longmapsto \varepsilon_\infty(x)=(w_j)_j
 \longmapsto\bigl(\Gamma_{A_j}(w_j)\bigr)_j
 \in\prod_j C_W(A_j),\\
C_W
&\longrightarrow W_{12}\xrightarrow{\ (D,\ell)\ }W_{24},\\
W_{12}&\xrightarrow{\ \operatorname{Aut}\ }M_{12},
\qquad W_{24}\xrightarrow{\ \operatorname{Aut}\ }M_{24},\\
C_W
&\longrightarrow N(A_2^{12})\longrightarrow\Lambda_{24}
 \longrightarrow V_{\Lambda_{24}}\longrightarrow V^\natural
 \longrightarrow\mathbb M\longrightarrow\text{Moonshine}.
\end{aligned}
\]
El peso y la distancia de Hamming pertenecen a la métrica de los códigos;
Hadamard pertenece a la transformación dodecafásica; Paley fija el calibre
de conferencia. Esta separación de funciones permite seguir cada flecha por
su operador, su dominio y su certificado, al tiempo que conserva la unidad
genealógica de la doble proyección.
